\section{A deployment-facing pilot: the cyber observer gate}

We also built a small observer-only fixture motivated by cyber-safety control.
The intended target was unsafe actionable assistance rather than security
topic. The nuisance labels were benign IT and defensive security. Matched
triples placed the same topic in each category, and the proposed evaluation
compared a text TF--IDF baseline, a linear activation observer, a
nuisance-residualized observer, a lifted observer, and a small MLP.

The fixture fails before it reaches representation analysis. With 15 test
examples in each class, the TF--IDF baseline has AUROC $1.0$, recalls all 15
unsafe examples, and makes no false positives on the 30 benign or defensive
examples. A linear activation observer also reaches AUROC $1.0$. The prompt
states intent in words, so the label is lexically exposed. The result cannot
establish that the residual stream contains a deployable unsafe-actionability
state.

The nuisance-residualized arm recalls $11/15$ unsafe examples and produces
false positives on $7/15$ benign-IT and $12/15$ defensive-security examples.
On this fixture, those counts do not diagnose residualization. Removing lexical
directions from a lexically determined target predictably removes target signal.
The measurement design confounds the target with stated intent before any
internal observer is tested.

This pilot is retained in the Phase 0--1 checkpoint to document the failed
fixture. It is not part of the main evidence. The final manuscript moves it to
an appendix and treats a contextual-authorization dataset, policy/rules
baselines, and cross-context holdouts as future work.
